% =====================================================================
%  COURS DE CHIMIE — FICHE 1
%  Les chiffres significatifs
%  Document généré en LaTeX avec figures TikZ tracées « from scratch ».
%  Compilation : pdflatex (2 passes pour la table des matières).
% =====================================================================
\documentclass[11pt,a4paper]{article}

% ---------- Encodage & langue ----------
\usepackage[utf8]{inputenc}
\usepackage[T1]{fontenc}
\usepackage{lmodern}
\usepackage[french]{babel}

% ---------- Mathématiques ----------
\usepackage{amsmath,amssymb,mathtools}
\usepackage{icomma}              % virgule décimale correcte en mode maths

% ---------- Unités & chimie ----------
\usepackage{siunitx}
\sisetup{output-decimal-marker={,},per-mode=symbol}
\usepackage[version=4]{mhchem}   % formules chimiques : \ce{NaCl}, \ce{NO2}...

% ---------- Couleurs, boîtes, graphismes ----------
\usepackage{xcolor}
\usepackage{colortbl}   % \rowcolor dans les tableaux
\usepackage{tikz}
\usepackage{pgfplots}
\usepackage[most]{tcolorbox}
\usepackage{enumitem}
\usepackage{array}
\usepackage{booktabs}
\usepackage{multicol}
\usepackage{fancyhdr}
\usepackage[hidelinks]{hyperref}

\usetikzlibrary{arrows.meta,positioning,calc,decorations.pathmorphing,
  decorations.markings,angles,quotes,patterns,backgrounds,
  shapes.geometric,shapes.misc,fadings,intersections,fit}
\pgfplotsset{compat=1.18}

% ---------- Géométrie de page ----------
\usepackage[a4paper,margin=2.2cm,headheight=26pt,headsep=10pt]{geometry}

% =====================================================================
%  PALETTE DE COULEURS  (mauve = couleur de la section Chimie du livre)
% =====================================================================
\definecolor{primary}{RGB}{150,98,162}      % mauve (couleur du livre — Chimie)
\definecolor{primarydark}{RGB}{92,52,110}
\definecolor{defcol}{RGB}{0,114,178}        % bleu  — définitions
\definecolor{thmcol}{RGB}{0,140,120}        % vert-bleu — théorèmes
\definecolor{formulacol}{RGB}{198,93,9}     % orange — formules
\definecolor{reglecol}{RGB}{150,98,162}     % mauve — règles
\definecolor{remarkcol}{RGB}{120,94,200}    % violet — remarques
\definecolor{attncol}{RGB}{200,40,55}       % rouge — pièges
\definecolor{excol}{RGB}{0,135,90}          % vert — exemples
\definecolor{methcol}{RGB}{90,90,100}       % gris — méthode
\definecolor{sigcol}{RGB}{0,90,200}         % bleu  — chiffres significatifs
\definecolor{nosigcol}{RGB}{210,30,40}      % rouge — chiffres non significatifs
\definecolor{certaincol}{RGB}{0,135,90}     % vert  — chiffres certains
\definecolor{doutecol}{RGB}{198,93,9}       % orange — chiffre incertain

% =====================================================================
%  STYLE COMMUN DES BOÎTES
% =====================================================================
\tcbset{
  boxbase/.style={enhanced,breakable,boxrule=0.9pt,arc=2.5pt,
     left=3mm,right=3mm,top=2.3mm,bottom=2.3mm,
     fonttitle=\bfseries,coltitle=white,
     toptitle=0.6mm,bottomtitle=0.6mm},
}

% ---- Définition (numérotée) ----
\newtcolorbox[auto counter,number within=section]{definition}[1][]{%
  boxbase,colback=defcol!5,colframe=defcol,
  title={\textsc{Définition}~\thetcbcounter\ \ifx\relax#1\relax\relax\else\textmd{--- \itshape #1}\fi}}

% ---- Théorème / Propriété (numéroté) ----
\newtcolorbox[auto counter,number within=section]{theoreme}[1][]{%
  boxbase,colback=thmcol!6,colframe=thmcol,
  title={\textsc{Propriété}~\thetcbcounter\ \ifx\relax#1\relax\relax\else\textmd{--- \itshape #1}\fi}}

% ---- Règle (numérotée) — spécifique aux chiffres significatifs ----
\newtcolorbox[auto counter,number within=section]{regle}[1][]{%
  boxbase,colback=reglecol!7,colframe=reglecol,
  title={\textsc{Règle}~\thetcbcounter\ \ifx\relax#1\relax\relax\else\textmd{--- \itshape #1}\fi}}

% ---- Formule (numérotée) ----
\newtcolorbox[auto counter,number within=section]{formule}[1][]{%
  boxbase,colback=formulacol!6,colframe=formulacol,
  title={\textsc{Formule}~\thetcbcounter\ \ifx\relax#1\relax\relax\else\textmd{--- \itshape #1}\fi}}

% ---- Remarque ----
\newtcolorbox{remarque}[1][]{%
  boxbase,colback=remarkcol!6,colframe=remarkcol,
  title={\textsc{Remarque}\ifx\relax#1\relax\relax\else\textmd{ : \itshape #1}\fi}}

% ---- Attention / piège ----
\newtcolorbox{attention}[1][]{%
  boxbase,colback=attncol!6,colframe=attncol,
  title={\textsc{Attention}\ifx\relax#1\relax\relax\else\textmd{ : \itshape #1}\fi}}

% ---- Exemple ----
\newtcolorbox{exemple}[1][]{%
  boxbase,colback=excol!5,colframe=excol,
  title={\textsc{Exemple}\ifx\relax#1\relax\relax\else\textmd{ : \itshape #1}\fi}}

% ---- Méthode ----
\newtcolorbox{methode}[1][]{%
  boxbase,colback=methcol!7,colframe=methcol,sharp corners,
  title={\textsc{Méthode}\ifx\relax#1\relax\relax\else\textmd{ : \itshape #1}\fi}}

% =====================================================================
%  ENVIRONNEMENT « où : »  (légende des grandeurs d'une formule)
% =====================================================================
\newenvironment{ou}{%
  \par\smallskip\noindent\textbf{où :}\nopagebreak
  \begin{itemize}[leftmargin=1.8em,topsep=2pt,itemsep=1.5pt,parsep=0pt]}%
  {\end{itemize}}

% =====================================================================
%  EN-TÊTES ET PIEDS DE PAGE
% =====================================================================
\pagestyle{fancy}
\fancyhf{}
\fancyhead[L]{\small\textcolor{primarydark}{\textbf{Fiche 1} — Les chiffres significatifs}}
\fancyhead[R]{\small\textcolor{primarydark}{\textsc{Chimie}}}
\fancyfoot[C]{\small\thepage}
\renewcommand{\headrulewidth}{0.8pt}
\renewcommand{\headrule}{\hbox to\headwidth{\color{primary}\leaders\hrule height \headrulewidth\hfill}}

% Titres de section en couleur
\usepackage{titlesec}
\titleformat{\section}{\Large\bfseries\color{primarydark}}{\thesection}{0.6em}{}
\titleformat{\subsection}{\large\bfseries\color{primary}}{\thesubsection}{0.6em}{}
\titleformat{\subsubsection}{\normalsize\bfseries\color{primary!80!black}}{\thesubsubsection}{0.6em}{}

% Raccourcis
\newcommand{\cs}{chiffre significatif}
\newcommand{\CS}{chiffres significatifs}

% =====================================================================
\begin{document}
\sloppy

% ---------------------------------------------------------------------
%  PAGE DE TITRE
% ---------------------------------------------------------------------
\begingroup
\thispagestyle{empty}
\centering
\vspace*{1.5cm}

\begin{tikzpicture}
\node[rectangle,rounded corners=8pt,fill=primary,text=white,
      minimum width=15.5cm,minimum height=2.6cm,align=center,inner sep=10pt]
  {{\Huge\bfseries Les chiffres significatifs}\\[6pt]
   {\large\bfseries Mesure, précision et incertitude en chimie}};
\end{tikzpicture}

\vspace{0.4cm}
{\large\color{primarydark}\textsc{Cours de chimie — Fiche n\textsuperscript{o}\,1}}

\vspace{0.9cm}

% --- Illustration de couverture : le nombre 0,0035070 décortiqué ---
\begin{tikzpicture}[scale=1.0]
  \node[font=\Large] (n) {$m=0,\textcolor{nosigcol}{\mathbf{00}}\textcolor{sigcol}{\mathbf{35070}}\ \si{kg}$};
  % accolade chiffres non significatifs (les deux zéros de tête)
  \draw[nosigcol,thick] (-1.62,0.35) -- (-1.62,0.62) -- (-1.18,0.62) -- (-1.18,0.35);
  \node[nosigcol,above,font=\scriptsize] at (-1.40,0.62) {zéros \emph{non} significatifs};
  % accolade chiffres significatifs
  \draw[sigcol,thick] (-0.95,-0.35) -- (-0.95,-0.62) -- (0.55,-0.62) -- (0.55,-0.35);
  \node[sigcol,below,font=\scriptsize] at (-0.20,-0.62) {5 chiffres significatifs};
\end{tikzpicture}

\vfill

\begin{tcolorbox}[boxbase,colback=primary!5,colframe=primary,width=15cm,
    title={\textsc{Objectifs pédagogiques de la fiche}}]
À l'issue de ce cours, l'étudiant·e sera capable de :
\begin{itemize}[leftmargin=1.6em,topsep=2pt,itemsep=1pt]
  \item comprendre pourquoi toute \emph{mesure} en chimie est entachée d'une \emph{incertitude} ;
  \item définir un \textbf{chiffre significatif} et distinguer chiffres \emph{certains} et chiffre \emph{incertain} ;
  \item appliquer les \textbf{règles de comptage} des \CS\ (cas délicat des zéros) ;
  \item maîtriser la \textbf{notation scientifique} et l'opération d'\textbf{arrondi} ;
  \item donner le bon nombre de \CS\ au résultat d'une \textbf{addition}, d'une \textbf{multiplication} et d'un \textbf{logarithme} (pH) ;
  \item distinguer \textbf{justesse} et \textbf{fidélité}, et exprimer une \textbf{incertitude} absolue et relative.
\end{itemize}
\end{tcolorbox}

\vspace{0.6cm}
{\small\color{black!60} Document de synthèse rédigé d'après la Fiche 1 du livre — figures originales tracées en TikZ.}
\par
\endgroup

% ---------------------------------------------------------------------
%  TABLE DES MATIÈRES
% ---------------------------------------------------------------------
\newpage
\renewcommand{\contentsname}{\textcolor{primarydark}{Table des matières}}
\tableofcontents
\thispagestyle{fancy}
\newpage

% =====================================================================
\section{Mesurer en chimie : grandeurs, nombres et incertitude}
% =====================================================================
La chimie est une science \textbf{quantitative} : peser un échantillon, lire un volume sur une
burette, relever une température, mesurer un pH\dots\ tout repose sur des \emph{mesures}. Or aucune
mesure n'est parfaite : tout instrument possède une \textbf{limite de précision}. Les
\textbf{\CS} sont l'outil qui permet d'écrire un nombre de façon à ce qu'il \emph{dise la vérité sur
sa propre précision} — ni plus, ni moins.

\begin{definition}[grandeur, valeur et unité]
Une \textbf{grandeur} est une propriété mesurable (une masse, un volume, une concentration). Le
résultat d'une mesure s'écrit toujours sous la forme d'un \textbf{nombre} suivi d'une \textbf{unité} :
\[ \underbrace{m}_{\text{grandeur}} \;=\; \underbrace{13{,}02}_{\text{nombre}}\ \underbrace{\si{\gram}}_{\text{unité}}. \]
Le nombre seul n'a aucun sens physique : \og $13{,}02$\fg{} ne veut rien dire, mais
\og $\qty{13{,}02}{\gram}$\fg{} désigne une masse précise.
\end{definition}

\begin{definition}[nombre mesuré et nombre exact]
On distingue deux natures de nombres :
\begin{itemize}[leftmargin=1.6em,topsep=2pt,itemsep=1.5pt]
  \item un \textbf{nombre mesuré} provient d'un instrument : il est toujours \emph{approché}, car
        l'instrument a une précision finie (ex.\ : $\qty{20{,}96}{\milli\litre}$ lu sur une burette) ;
  \item un \textbf{nombre exact} provient d'un comptage ou d'une définition : il est \emph{parfaitement}
        connu (ex.\ : il y a \emph{exactement} $2$ atomes d'hydrogène dans \ce{H2O}, ou
        \emph{exactement} $\qty{1000}{\gram}$ dans $\qty{1}{\kilo\gram}$ par définition).
\end{itemize}
\end{definition}

\begin{remarque}
La règle des \CS\ ne s'applique qu'aux \emph{nombres mesurés}. Un nombre exact est considéré comme
ayant une \emph{infinité} de \CS\ : il ne limite donc jamais la précision d'un calcul (on y revient
au \S\ref{sec:exacts}).
\end{remarque}

\subsection{L'exemple de la burette}

Le contexte le plus parlant est la lecture d'un volume sur une \textbf{burette} graduée, instrument
de base du laboratoire de chimie (titrages).

\begin{center}
\begin{tikzpicture}[scale=1.0]
  % corps de la burette
  \draw[thick,rounded corners=2pt] (0,0) rectangle (1.1,6);
  \draw[thick] (0.45,0) -- (0.45,-0.7);          % robinet/pointe
  \draw[thick] (0.65,0) -- (0.65,-0.7);
  \fill[primary!20] (0,0) rectangle (1.1,2.55);  % liquide
  % graduations principales 20 et 21
  \foreach \y/\lab in {1.2/21,4.8/20}{
     \draw[thick] (1.1,\y) -- (1.5,\y);
     \node[right,font=\small] at (1.5,\y) {\lab};}
  % sous-graduations (1/10 de mL)
  \foreach \y in {1.56,1.92,2.28,2.64,3.0,3.36,3.72,4.08,4.44}{
     \draw (1.1,\y) -- (1.35,\y);}
  % ménisque
  \draw[primary,thick] (0,2.55) .. controls (0.4,2.4) and (0.7,2.4) .. (1.1,2.55);
  \draw[->,>=Stealth,nosigcol,thick] (3.0,2.55) -- (1.2,2.55);
  \node[right,nosigcol,font=\small] at (3.0,2.55) {ménisque $\approx \qty{20{,}96}{\milli\litre}$};
  % zone d'incertitude
  \draw[densely dashed,doutecol] (-0.3,2.16) -- (1.1,2.16);
  \draw[densely dashed,doutecol] (-0.3,2.94) -- (1.1,2.94);
  \node[left,doutecol,font=\scriptsize,align=right] at (-0.35,2.55)
       {entre\\ \num{20,9} et \num{21,0}};
\end{tikzpicture}
\end{center}

\noindent Si la plus petite graduation vaut $\qty{0{,}1}{\milli\litre}$, on voit que le ménisque
se situe \emph{entre} $\qty{20{,}9}{}$ et $\qty{21{,}0}{\milli\litre}$. À l'\oe il nu, on
\emph{estime} la position à $\qty{20{,}96}{\milli\litre}$.

\begin{itemize}[leftmargin=1.6em,topsep=3pt]
  \item Les chiffres \textbf{2}, \textbf{0} et \textbf{9} sont \textcolor{certaincol}{\textbf{certains}} :
        la graduation nous les garantit.
  \item Le dernier chiffre \textbf{6} est \textcolor{doutecol}{\textbf{incertain}} (estimé) : un autre
        observateur aurait pu lire $5$ ou $7$. Il n'y a pas de graduation au centième de \si{\milli\litre}.
\end{itemize}

\begin{theoreme}[plus de \CS\ $\Rightarrow$ plus de précision]
En général, \emph{plus une mesure comporte de \CS, plus elle est précise}. Le nombre de \CS\ est une
\og carte d'identité\fg{} de la qualité de la mesure : il indique jusqu'où l'instrument est fiable.
Le volume $V=\qty{20{,}96}{\milli\litre}$ comporte \textbf{quatre} \CS.
\end{theoreme}

% =====================================================================
\section{La notion de chiffre significatif}
% =====================================================================

\begin{definition}[chiffre significatif]
Les \textbf{\CS} d'un nombre mesuré sont \emph{tous les chiffres connus avec certitude}, \emph{plus le
premier chiffre incertain} (le dernier écrit, celui qui est estimé). Ce sont les chiffres qui
\emph{portent une information sur la mesure}.
\end{definition}

\begin{center}
\begin{tikzpicture}[font=\small]
  \node[font=\Large] (n) at (0,0) {$V = \textcolor{certaincol}{20{,}9}\textcolor{doutecol}{6}\ \si{\milli\litre}$};
  \draw[certaincol,thick,->,>=Stealth] (-0.85,-0.55) -- (-0.85,-0.15);
  \node[certaincol,below,align=center,font=\scriptsize] at (-0.85,-0.55)
       {3 chiffres \textbf{certains}\\ (lus sur la graduation)};
  \draw[doutecol,thick,->,>=Stealth] (0.55,0.55) -- (0.55,0.18);
  \node[doutecol,above,align=center,font=\scriptsize] at (0.55,0.55)
       {1 chiffre \textbf{incertain}\\ (estimé à l'\oe il)};
  \node[primarydark,right,font=\footnotesize] at (1.9,0)
       {$\Rightarrow\ 3+1 = \mathbf{4}$ chiffres significatifs};
\end{tikzpicture}
\end{center}

\begin{theoreme}[les chiffres non nuls sont toujours significatifs]
\textbf{Tout chiffre non nul} ($1,2,3,4,5,6,7,8,9$) est \emph{toujours} significatif, quelle que soit
sa position dans le nombre. La seule difficulté de comptage concerne le chiffre \textbf{0} (\S\ref{sec:zeros}).
\end{theoreme}

\begin{exemple}[comptage immédiat]
\begin{itemize}[leftmargin=1.6em,topsep=2pt]
  \item $\num{837}$ : trois chiffres non nuls $\Rightarrow$ \textbf{3} \CS.
  \item $\num{42,5}$ : trois chiffres non nuls $\Rightarrow$ \textbf{3} \CS.
  \item $\num{6,1234}$ : cinq chiffres non nuls $\Rightarrow$ \textbf{5} \CS.
\end{itemize}
\end{exemple}

\begin{attention}[chiffre significatif $\neq$ nombre de décimales]
Ne confondez pas \emph{\CS} et \emph{décimales} (chiffres après la virgule). $\num{42,5}$ a une seule
décimale mais \emph{trois} \CS ; $\num{0,007}$ a trois décimales mais un \emph{seul} \cs. Le nombre de
\CS\ se compte à partir du \textbf{premier chiffre non nul}, en lisant de gauche à droite.
\end{attention}

% =====================================================================
\section{Compter les chiffres significatifs : le cas des zéros}
\label{sec:zeros}
% =====================================================================
Le chiffre $0$ joue deux rôles très différents : tantôt il \emph{cale} la virgule (il ne fait que
positionner l'ordre de grandeur), tantôt il \emph{résulte d'une mesure} (il porte alors une
information). Trois règles règlent tous les cas.

\begin{regle}[zéros de tête — \emph{jamais} significatifs]
Les zéros situés \textbf{à gauche} du premier chiffre non nul ne sont \emph{jamais} significatifs :
ils ne servent qu'à placer la virgule (l'ordre de grandeur). On les appelle \emph{zéros de cadrage}.
\end{regle}

\begin{regle}[zéros intercalés — \emph{toujours} significatifs]
Les zéros situés \textbf{entre} deux chiffres non nuls (zéros \emph{captifs}) sont \emph{toujours}
significatifs.
\end{regle}

\begin{regle}[zéros de queue après la virgule — \emph{toujours} significatifs]
Les zéros situés \textbf{à droite}, \emph{après la virgule} (et après le dernier chiffre non nul),
sont \emph{toujours} significatifs : on ne les écrirait pas s'ils n'étaient pas le fruit d'une mesure.
\end{regle}

\subsection{La figure-clé : $m = \mathbf{0{,}0035070}$ kg}

L'exemple suivant, tiré du livre, contient \emph{tous} les cas de figure à la fois.

\begin{center}
\begin{tikzpicture}[scale=1.0]
  \node[font=\LARGE] (n) {$m\;=\;0,00\textcolor{nosigcol}{\mathbf{3}}\textcolor{nosigcol}{\mathbf{5}}%
        \textcolor{sigcol}{\mathbf{0}}\textcolor{nosigcol}{\mathbf{7}}\textcolor{sigcol}{\mathbf{0}}\ \si{kg}$};
\end{tikzpicture}

\medskip
\begin{tikzpicture}[font=\small]
  \node[font=\LARGE] (n) at (0,0)
     {$0,\textcolor{nosigcol}{00}\,\textbf{3}\,\textbf{5}\,\textcolor{sigcol}{\textbf 0}\,\textbf{7}\,\textcolor{sigcol}{\textbf 0}$};
  % zéros non significatifs (tête)
  \draw[nosigcol,thick,->,>=Stealth] (-1.05,0.75) -- (-1.05,0.25);
  \node[nosigcol,above,align=center,font=\scriptsize] at (-1.05,0.75)
       {les \textbf{zéros de tête}\\ ne sont \emph{pas} significatifs};
  % zéros significatifs (queue / captif)
  \draw[sigcol,thick,->,>=Stealth] (0.30,-0.78) -- (0.30,-0.25);
  \draw[sigcol,thick,->,>=Stealth] (1.18,-0.78) -- (1.18,-0.25);
  \node[sigcol,below,align=center,font=\scriptsize] at (0.74,-0.85)
       {ces \textbf{zéros} (captif et de queue)\\ \emph{sont} significatifs};
\end{tikzpicture}
\end{center}

\noindent\textbf{Lecture chiffre par chiffre} (de gauche à droite) :
\begin{itemize}[leftmargin=1.6em,topsep=3pt,itemsep=2pt]
  \item $\textcolor{nosigcol}{0,00}$ : trois zéros de tête $\rightarrow$ \emph{cadrage}, \textbf{non} significatifs ;
  \item $\mathbf{3}$, $\mathbf{5}$ : non nuls $\rightarrow$ significatifs ;
  \item $\textcolor{sigcol}{\mathbf{0}}$ : zéro \emph{captif} (entre le $5$ et le $7$) $\rightarrow$ significatif ;
  \item $\mathbf{7}$ : non nul $\rightarrow$ significatif ;
  \item $\textcolor{sigcol}{\mathbf{0}}$ : zéro \emph{de queue} après la virgule $\rightarrow$ significatif.
\end{itemize}
\noindent Bilan : $3,\,5,\,0,\,7,\,0$ sont significatifs, soit $\boxed{\mathbf 5}$ \CS.

\begin{theoreme}[le déplacement de la virgule ne change rien]
\label{thm:virgule}
Le \textbf{déplacement de la virgule} (donc le changement d'unité ou la notation scientifique)
n'a \emph{aucun} impact sur la précision : il ne modifie jamais le nombre de \CS. Seuls changent les
\emph{zéros de cadrage}, qui ne comptent pas.
\end{theoreme}

\begin{exemple}[un même volume, quatre \CS\ dans toutes les unités]
Le volume $V=\qty{20{,}96}{\milli\litre}$ possède $4$ \CS. Écrivons-le autrement :
\[ V=\underbrace{\num{20,96}}_{4\ \text{CS}}\ \si{\milli\litre}
    =\underbrace{\num{0,02096}}_{4\ \text{CS}}\ \si{\litre}
    =\underbrace{\num{2,096e-2}}_{4\ \text{CS}}\ \si{\litre}. \]
Dans $\num{0,02096}\,\si{\litre}$, les deux zéros de tête ($0{,}0$) ne sont \textbf{pas} significatifs ;
ils disparaissent en notation scientifique $\num{2,096e-2}$. Le nombre de \CS\ reste $4$ partout :
la \emph{précision} de la mesure n'a pas changé, seule l'écriture a changé.
\end{exemple}

\subsection{Le piège des zéros de queue sans virgule}

Un cas reste \emph{ambigu} : les zéros de queue d'un nombre \textbf{entier} (sans virgule).

\begin{attention}[combien de \CS\ dans \og $\num{100}$\fg{} ?]
Écrit tel quel, $\num{100}$ est \emph{ambigu} : on ne sait pas si les zéros résultent d'une mesure
ou d'un simple cadrage. Le nombre peut avoir $1$, $2$ ou $3$ \CS. La \textbf{notation scientifique}
lève toute ambiguïté :
\begin{center}\renewcommand{\arraystretch}{1.3}
\begin{tabular}{c c c}
\toprule
\textbf{Écriture} & \textbf{Nombre de \CS} & \textbf{Signification}\\
\midrule
$\num{1e2}$    & $1$ & mesure au $\pm\,\num{100}$ près\\
$\num{1,0e2}$  & $2$ & mesure à $\pm\,\num{10}$ près\\
$\num{1,00e2}$ & $3$ & mesure à $\pm\,\num{1}$ près\\
\bottomrule
\end{tabular}
\end{center}
C'est pourquoi, en chimie, on privilégie systématiquement la notation scientifique dès qu'un doute
est possible.
\end{attention}

\begin{exemple}[reconnaître les zéros significatifs]
\begin{center}\renewcommand{\arraystretch}{1.35}
\begin{tabular}{>{\centering\arraybackslash}m{2.6cm} c >{\raggedright\arraybackslash}m{7.4cm}}
\toprule
\textbf{Nombre} & \textbf{\CS} & \textbf{Justification}\\
\midrule
$\num{15,200}$    & $5$ & les deux zéros de queue (après virgule) sont significatifs\\
$\num{0,0100}$    & $3$ & zéros de tête non sig.\ ; les deux zéros de queue le sont\\
$\num{0,0035070}$ & $5$ & zéros de tête non sig.\ ; le captif et celui de queue le sont\\
$\num{2030}$      & $3$ (ambigu) & le zéro captif est sig.\ ; le zéro de queue est ambigu\\
$\num{4,00e3}$    & $3$ & notation scientifique : tout ce qui est écrit est significatif\\
\bottomrule
\end{tabular}
\end{center}
\end{exemple}

% =====================================================================
\section{Nombres exacts : une infinité de chiffres significatifs}
\label{sec:exacts}
% =====================================================================
Certains nombres ne proviennent d'aucune mesure : ils sont \emph{exacts} et ne limitent donc jamais
la précision.

\begin{definition}[nombre exact]
Un \textbf{nombre exact} est connu sans aucune incertitude. Deux origines :
\begin{itemize}[leftmargin=1.6em,topsep=2pt,itemsep=1.5pt]
  \item un \textbf{comptage} d'objets discrets : $3$ béchers, $2$ atomes de \ce{H} dans \ce{H2O},
        $6$ faces d'un dé ;
  \item une \textbf{définition} ou une relation conventionnelle : $\qty{1}{\kilo\gram}=\qty{1000}{\gram}$
        \emph{exactement}, $\qty{1}{\hour}=\qty{3600}{\second}$ \emph{exactement}, le facteur $2$ dans
        un périmètre $P=2\pi r$.
\end{itemize}
On lui attribue une \textbf{infinité} de \CS ; il n'intervient donc jamais dans la limitation du
résultat.
\end{definition}

\begin{exemple}[un facteur de conversion ne dégrade pas la précision]
On convertit une masse $m=\qty{13{,}02}{\gram}$ (4 \CS) en kilogrammes. Le facteur
$\qty{1}{\kilo\gram}=\qty{1000}{\gram}$ est \emph{exact} :
\[ m=\frac{\num{13,02}}{\num{1000}}\ \si{\kilo\gram}=\qty{0,01302}{\kilo\gram}. \]
Le résultat conserve ses $4$ \CS : le $\num{1000}$ exact n'a rien limité. C'est la propriété du
théorème \ref{thm:virgule} : on n'a fait que déplacer la virgule.
\end{exemple}

\begin{remarque}
Les constantes mathématiques comme $\pi$ ou $\sqrt{2}$ sont, elles aussi, connues avec autant de
décimales qu'on veut : on les prend avec \emph{au moins} autant de \CS\ que la donnée la plus précise,
pour ne pas être le maillon faible du calcul.
\end{remarque}

% =====================================================================
\section{La notation scientifique}
% =====================================================================
La notation scientifique est l'écriture \emph{reine} de la chimie : elle gère les très grands nombres
(constante d'Avogadro $\approx \qty{6,022e23}{\per\mole}$) comme les très petits
(concentration en ions \ce{H+} d'une solution neutre $=\qty{1,0e-7}{\mol\per\litre}$), tout en
affichant sans ambiguïté le nombre de \CS.

\begin{definition}[notation scientifique]
Un nombre est en \textbf{notation scientifique} lorsqu'il s'écrit
\[ \boxed{\;a\times 10^{\,n}\;}\qquad\text{avec}\quad 1\le |a| < 10\ \text{et}\ n\in\mathbb{Z}. \]
\begin{ou}
  \item $a$ : la \textbf{mantisse}, un nombre dont la partie entière est \emph{un seul} chiffre non nul ; elle porte \emph{tous} les \CS ;
  \item $10^{\,n}$ : la \textbf{puissance de dix}, qui fixe l'\emph{ordre de grandeur} (n'apporte aucun \cs) ;
  \item $n$ : un entier relatif (positif pour les grands nombres, négatif pour les petits).
\end{ou}
\end{definition}

\begin{methode}[mettre un nombre en notation scientifique]
\begin{enumerate}[leftmargin=1.8em,topsep=2pt,itemsep=1.5pt]
  \item placer la virgule \emph{juste après le premier chiffre non nul} ;
  \item compter de combien de rangs la virgule a été déplacée : c'est $|n|$ ;
  \item $n>0$ si le nombre de départ était \og grand\fg{} ($\ge 10$), $n<0$ s'il était \og petit\fg{} ($<1$) ;
  \item ne garder dans la mantisse que les \CS\ (on supprime les zéros de cadrage).
\end{enumerate}
\end{methode}

\begin{exemple}[aller-retour avec la notation scientifique]
\begin{align*}
 \num{20960}      &= \num{2,0960e4} &&(\text{virgule reculée de 4 rangs})\\
 \num{0,02096}    &= \num{2,096e-2} &&(\text{virgule avancée de 2 rangs})\\
 \num{6,022e23}   &= \num{602200000000000000000000} &&(\text{constante d'Avogadro})\\
 \num{1,0e-7}     &= \num{0,00000010} &&(\ce{[H+]} \text{d'une solution neutre})
\end{align*}
\end{exemple}

\begin{exemple}[une même pesée à 4 \CS\ dans cinq unités]
Une balance affiche $m=\qty{13{,}02}{\gram}$ ($4$ \CS). La même masse s'écrit, sans jamais perdre
ni gagner de \cs\ (théorème \ref{thm:virgule}) :
\begin{center}\renewcommand{\arraystretch}{1.3}
\begin{tabular}{l l c}
\toprule
\textbf{Unité} & \textbf{Valeur} & \textbf{\CS}\\
\midrule
gramme       & $\qty{13{,}02}{\gram}$            & $4$\\
milligramme  & $\num{1,302e4}\ \si{\milli\gram}$ \;=\; $\qty{13020}{\milli\gram}$ & $4$\\
kilogramme   & $\qty{0,01302}{\kilo\gram}$       & $4$\\
tonne        & $\qty{0,00001302}{\tonne}$ \;=\; $\num{1,302e-5}\ \si{\tonne}$ & $4$\\
\bottomrule
\end{tabular}
\end{center}
\textbf{Interprétation.} Écrire $\qty{13{,}02}{\gram}$ est, par convention, une façon abrégée de dire
\[ m = \qty{13{,}02\pm0{,}01}{\gram}, \]
c'est-à-dire que la vraie valeur est comprise entre $\qty{13{,}01}{}$ et $\qty{13{,}03}{\gram}$.
Le dernier chiffre écrit (le $2$) est le chiffre \emph{incertain}.
\end{exemple}

\begin{exemple}[la même pesée annoncée à 3 \CS]
Si la balance n'est précise qu'au dixième de gramme, la pesée vaut $\qty{13{,}0}{\gram}$ ($3$ \CS,
le zéro de queue étant significatif !). Dans les autres unités :
\[ \qty{13{,}0}{\gram}=\num{1,30e4}\ \si{\milli\gram}=\qty{0,0130}{\kilo\gram}
   =\qty{0,0000130}{\tonne}=\num{1,30e-5}\ \si{\tonne}. \]
On \emph{n'écrit pas} $\qty{13}{\gram}$ : cela ne ferait que $2$ \CS\ et \emph{perdrait} l'information
\og précis au dixième\fg{}. Le zéro de queue est ici porteur d'information.
\end{exemple}

\begin{remarque}[notation \emph{ingénieur}]
Variante de la notation scientifique où l'exposant $n$ est un \emph{multiple de 3}, ce qui colle aux
préfixes du Système International (kilo $10^3$, milli $10^{-3}$, micro $10^{-6}$, nano $10^{-9}$\dots).
Exemple : $\num{2,096e-2}\,\si{\litre}=\num{20,96e-3}\,\si{\litre}=\qty{20,96}{\milli\litre}$.
\end{remarque}

% =====================================================================
\section{Arrondir un résultat}
% =====================================================================
Réduire un nombre au bon nombre de \CS\ exige de savoir \textbf{arrondir} proprement.

\begin{regle}[arrondi au rang voulu]
Pour arrondir, on regarde le \textbf{premier chiffre supprimé} (celui juste à droite du dernier
chiffre conservé) :
\begin{itemize}[leftmargin=1.6em,topsep=2pt,itemsep=1.5pt]
  \item s'il vaut $0,1,2,3,4$ : on \textbf{arrondit par défaut} (on ne change pas le dernier chiffre gardé) ;
  \item s'il vaut $5,6,7,8,9$ : on \textbf{arrondit par excès} (on ajoute $1$ au dernier chiffre gardé).
\end{itemize}
\end{regle}

\begin{center}
\begin{tikzpicture}[font=\small]
  \draw[->,>=Stealth,thick,black!70] (0,0) -- (10.4,0);
  \foreach \x/\d in {0/0,1.15/1,2.3/2,3.45/3,4.6/4,5.75/5,6.9/6,8.05/7,9.2/8,10.35/9}{
     \draw (\x,0.12) -- (\x,-0.12);
     \node[below,font=\footnotesize] at (\x,-0.14) {\d};}
  % zone défaut
  \draw[excol,line width=2pt] (0,0.32) -- (4.6,0.32);
  \node[excol,above,font=\footnotesize] at (2.3,0.34) {arrondi \textbf{par défaut} (0–4)};
  % zone excès
  \draw[attncol,line width=2pt] (5.75,0.32) -- (10.35,0.32);
  \node[attncol,above,font=\footnotesize] at (8.05,0.34) {arrondi \textbf{par excès} (5–9)};
  \draw[densely dashed,black!50] (5.175,-0.55) -- (5.175,0.7);
\end{tikzpicture}
\end{center}

\begin{exemple}[arrondir à un nombre de \CS\ donné]
\begin{center}\renewcommand{\arraystretch}{1.35}
\begin{tabular}{c c c c}
\toprule
\textbf{Nombre} & \textbf{À \dots \CS} & \textbf{Chiffre supprimé} & \textbf{Arrondi}\\
\midrule
$\num{3,148}$  & $3$ & $8\ (\ge5)$ & $\num{3,15}$ (par excès)\\
$\num{3,142}$  & $3$ & $2\ (<5)$   & $\num{3,14}$ (par défaut)\\
$\num{0,06237}$ & $2$ & $3\ (<5)$  & $\num{0,062}$\\
$\num{124,7}$  & $3$ & $7\ (\ge5)$ & $\num{125}$\\
$\num{2,996}$  & $3$ & $6\ (\ge5)$ & $\num{3,00}$ (retenue en cascade)\\
\bottomrule
\end{tabular}
\end{center}
La dernière ligne illustre une \textbf{retenue} : $\num{2,996}\to\num{3,00}$. Garder le \og $00$\fg{}
est obligatoire pour conserver $3$ \CS.
\end{exemple}

\begin{exemple}[le volume d'un cube — un calcul, puis un arrondi]
On veut le volume d'un cube d'arête $a=\qty{2{,}5}{\centi\metre}$, exprimé \emph{à 4 \CS}.
\[ V=a^3=\num{2,5}\times\num{2,5}\times\num{2,5}=\qty{15,625}{\centi\metre\cubed}. \]
Pour l'écrire à $4$ \CS, on supprime le chiffre $5$ (rang des millièmes) : comme $5\ge5$, on arrondit
\emph{par excès} :
\[ \boxed{V \approx \qty{15{,}63}{\centi\metre\cubed}}. \]
\begin{remarque}
Stricto sensu, l'arête $\num{2,5}$ n'ayant que $2$ \CS, la \emph{règle de multiplication}
(\S\ref{sec:opmult}) imposerait un résultat à $2$ \CS, soit $\qty{16}{\centi\metre\cubed}$. L'énoncé
demandant explicitement $4$ \CS, on se contente ici d'un exercice d'\emph{arrondi}. Retenez bien la
différence entre \og arrondir à $n$ \CS\fg{} (mécanique) et \og donner le bon nombre de \CS\fg{}
(qui dépend des données, cf.\ section suivante).
\end{remarque}
\end{exemple}

\begin{attention}[n'arrondir qu'à la \emph{fin}]
Dans un calcul à plusieurs étapes, on \textbf{garde des chiffres supplémentaires} (chiffres
\emph{de garde}) à chaque étape intermédiaire et on n'arrondit qu'au \textbf{résultat final}. Arrondir
trop tôt propage et amplifie l'erreur (\og erreur d'arrondi cumulée\fg{}).
\end{attention}

% =====================================================================
\section{Les opérations avec chiffres significatifs}
% =====================================================================
Le résultat d'un calcul ne peut pas être \emph{plus précis} que les mesures dont il est issu. Deux
règles distinctes s'appliquent : l'une pour l'addition/soustraction, l'autre pour la
multiplication/division.

\subsection{Addition et soustraction : on compte les \emph{décimales}}

\begin{regle}[addition / soustraction]
Le résultat d'une \textbf{addition} ou d'une \textbf{soustraction} s'arrondit au \emph{même nombre de
décimales} (chiffres après la virgule) que la donnée qui en a le \textbf{moins}.
\end{regle}

\noindent\textbf{Pourquoi les décimales et non les \CS ?} Parce qu'une somme se fait \emph{rang par
rang} : le rang le moins bien connu d'un des termes contamine ce même rang dans le résultat. C'est la
position de la dernière décimale fiable qui compte, pas le nombre total de chiffres.

\begin{exemple}[addition de trois volumes]
On additionne trois volumes lus sur des instruments de précisions différentes :
\[
\begin{array}{r@{}l}
  123{,}&25\quad(\text{2 décimales})\\
   46{,}&0\phantom{0}\quad(\text{\textbf{1} décimale} \leftarrow \text{le moins précis})\\
   86{,}&257\quad(\text{3 décimales})\\
  \hline
  255{,}&507
\end{array}
\]
La donnée la moins précise ($\num{46,0}$) n'a qu'\textbf{une} décimale ; on arrondit donc le résultat
à une décimale :
\[ \boxed{\num{123,25}+\num{46,0}+\num{86,257} \approx \num{255,5}}. \]
Inutile d'écrire $\num{255,507}$ : les chiffres $0$ et $7$ sont illusoires, car $\num{46,0}$ ignore
déjà tout ce qui se passe au-delà du dixième.
\end{exemple}

\begin{exemple}[soustraction — perte de \CS]
\[ \num{8,11} - \num{8,0} = \num{0,11}\ \xrightarrow{\text{1 décimale}}\ \num{0,1}. \]
Les deux nombres avaient $3$ et $2$ \CS, mais la différence n'en a plus qu'\textbf{un} ! La
soustraction de deux nombres voisins est une \emph{catastrophe} pour la précision (\og perte de
chiffres significatifs\fg{}) : à éviter dès qu'on le peut en chimie analytique.
\end{exemple}

\subsection{Multiplication et division : on compte les \CS}
\label{sec:opmult}

\begin{regle}[multiplication / division]
Le résultat d'une \textbf{multiplication} ou d'une \textbf{division} s'arrondit au \emph{même nombre de
\CS} que la donnée qui en a le \textbf{moins}.
\end{regle}

\begin{exemple}[concentration massique de \ce{NO2} dans l'air]
Une chimiste trouve $\qty{0,0361}{\milli\gram}$ de dioxyde d'azote \ce{NO2} dans un échantillon de
$\qty{12,0}{\litre}$ d'air. On cherche la concentration $[\ce{NO2}]$ en \si{\gram\per\cubic\metre}.
\[ [\ce{NO2}]=\frac{m}{V}=\frac{\qty{0,0361}{\milli\gram}}{\qty{12,0}{\litre}}
   =\frac{\num{0,0361e-3}\ \si{\gram}}{\num{12,0e-3}\ \si{\cubic\metre}}. \]
\textbf{Étape 1 — le calcul brut :}
\[ [\ce{NO2}]=\frac{\num{0,0361e-3}}{\num{12,0e-3}}\ \si{\gram\per\cubic\metre}
   =\num{0,00300833}\ \si{\gram\per\cubic\metre}=\num{3,00833e-3}\ \si{\gram\per\cubic\metre}. \]
\textbf{Étape 2 — combien de \CS ?} La donnée la moins riche est $\qty{12,0}{\litre}$, avec $3$ \CS
(autant que $\num{0,0361}$). On arrondit donc à $3$ \CS :
\[ \boxed{[\ce{NO2}]\approx\num{3,01e-3}\ \si{\gram\per\cubic\metre}}. \]
\textbf{Note d'unité.} On a utilisé $\qty{1}{\milli\gram}=\num{e-3}\,\si{\gram}$ et
$\qty{1}{\litre}=\num{e-3}\,\si{\cubic\metre}$ (facteurs \emph{exacts}). La concentration d'un gaz
dans l'air s'exprime couramment en \si{\gram\per\cubic\metre}.
\end{exemple}

\begin{exemple}[énergie de combustion par kilogramme]
La combustion de $\qty{1}{\gram}$ d'essence libère $\num{4,307e4}\ \si{\joule}$ ($4$ \CS). Quelle est
l'énergie libérée \emph{par kilogramme}, à $3$ \CS ?

Comme $\qty{1}{\gram}=\num{e-3}\,\si{\kilo\gram}$ (facteur exact) :
\[ \frac{\num{4,307e4}\ \si{\joule}}{\qty{1}{\gram}}
   =\frac{\num{4,307e4}\ \si{\joule}}{\num{e-3}\ \si{\kilo\gram}}
   =\num{4,307e7}\ \si{\joule\per\kilo\gram}. \]
Arrondi à $3$ \CS\ (chiffre supprimé $7\ge5$, par excès) :
\[ \boxed{\approx\num{4,31e7}\ \si{\joule\per\kilo\gram}}. \]
\end{exemple}

\begin{exemple}[solubilité — une règle de trois]
On peut dissoudre au maximum $\qty{36}{\gram}$ de chlorure de sodium \ce{NaCl} dans $\qty{100}{\gram}$
d'eau à $\qty{25}{\celsius}$. Quelle masse maximale de \ce{NaCl} se dissout dans $\qty{93,233}{\gram}$
d'eau à la même température ? (résultat en \si{\kilo\gram}, à $4$ \CS).

\textbf{Règle de trois.}
\[
\begin{array}{r c l}
 \qty{36}{\gram}\ \text{de \ce{NaCl}} &\longleftrightarrow& \qty{100}{\gram}\ \text{d'eau}\\[2pt]
 x\ \text{g de \ce{NaCl}}             &\longleftrightarrow& \qty{93,233}{\gram}\ \text{d'eau}
\end{array}
\qquad\Longrightarrow\qquad
x=\frac{\num{36}\times\num{93,233}}{\num{100}}.
\]
\textbf{Calcul brut :}
\[ x=\frac{\num{36}\times\num{93,233}}{\num{100}}=\frac{\num{3356,388}}{\num{100}}=\qty{33,56388}{\gram}. \]
En kilogrammes et à $4$ \CS :
\[ \boxed{x\approx\qty{0,03356}{\kilo\gram}=\num{3,356e-2}\ \si{\kilo\gram}}. \]
\begin{remarque}
L'énoncé impose $4$ \CS, on s'y plie. En toute rigueur cependant, $\qty{36}{\gram}$ n'a que $2$ \CS :
la règle de multiplication limiterait alors le résultat à $2$ \CS\ ($\approx\qty{34}{\gram}$). Cet
écart illustre l'importance d'écrire les données avec leur \emph{vraie} précision
($\qty{36{,}00}{\gram}$ si elle est connue à $4$ \CS).
\end{remarque}
\end{exemple}

\subsection{Calculs en chaîne et cas mixtes}

\begin{methode}[un calcul mêlant $+$ et $\times$]
\begin{enumerate}[leftmargin=1.8em,topsep=2pt,itemsep=1.5pt]
  \item effectuer le calcul \emph{en entier} sans arrondir, en gardant des chiffres de garde ;
  \item appliquer la règle adaptée à \emph{chaque} type d'opération, dans l'ordre des étapes ;
  \item arrondir \emph{une seule fois}, à la fin, au nombre de \CS\ (ou de décimales) imposé par
        l'étape limitante.
\end{enumerate}
\end{methode}

\begin{exemple}[masse molaire et nombre de moles]
La masse molaire de l'eau se calcule à partir de masses molaires atomiques
$M(\ce{H})=\qty{1,008}{\gram\per\mole}$ et $M(\ce{O})=\qty{16,00}{\gram\per\mole}$. Les indices ($2$ et
$1$) sont des nombres \emph{exacts} (comptage d'atomes) :
\[ M(\ce{H2O}) = 2\times\num{1,008} + \num{16,00}
   = \num{2,016} + \num{16,00} = \qty{18,016}{\gram\per\mole}. \]
C'est une \emph{addition} : on garde le nombre de décimales du moins précis ($\num{16,00}$, deux
décimales) $\Rightarrow M(\ce{H2O})\approx\qty{18,02}{\gram\per\mole}$.

\medskip
Combien de moles dans $m=\qty{9,00}{\gram}$ d'eau ? On \emph{divise} :
\[ n=\frac{m}{M}=\frac{\qty{9,00}{\gram}}{\qty{18,02}{\gram\per\mole}}=\num{0,4994}\ldots\ \si{\mole}. \]
Donnée la moins riche : $\qty{9,00}{\gram}$ ($3$ \CS) $\Rightarrow$
$\boxed{n\approx\qty{0,499}{\mole}}$.
\end{exemple}

\subsection{Logarithmes et pH : un cas propre à la chimie}

Le \textbf{pH} est défini par un logarithme décimal ; la règle des \CS\ y prend une forme
particulière, indispensable en chimie des solutions.

\begin{formule}[définition du pH]
\[ \boxed{\;\mathrm{pH}=-\log_{10}\bigl([\ce{H3O+}]\bigr)\;} \]
\begin{ou}
  \item $\mathrm{pH}$ : grandeur \emph{sans unité} caractérisant l'acidité d'une solution ;
  \item $[\ce{H3O+}]$ : concentration en ions oxonium, en \si{\mol\per\litre} ;
  \item $\log_{10}$ : logarithme décimal (fonction réciproque de $10^{x}$).
\end{ou}
\end{formule}

\begin{regle}[\CS\ d'un logarithme]
Dans un $\log$, seuls les chiffres \emph{après la virgule} (la \og mantisse\fg{} décimale) sont
significatifs ; la \emph{partie entière} ne fait que rappeler la puissance de dix. Donc :
\[ \text{nombre de \emph{décimales} du pH} \;=\; \text{nombre de \CS\ de } [\ce{H3O+}]. \]
\end{regle}

\begin{exemple}[calculer un pH avec le bon nombre de décimales]
Soit $[\ce{H3O+}]=\num{2,5e-4}\ \si{\mol\per\litre}$ : la concentration a \textbf{2} \CS.
\[ \mathrm{pH}=-\log(\num{2,5e-4})=-(\log\num{2,5}+\log\num{e-4})
   =-(0{,}39794\ldots-4)=3{,}60205\ldots \]
La concentration ayant $2$ \CS, on garde \textbf{2 décimales} :
\[ \boxed{\mathrm{pH}\approx 3{,}60}. \]
Les deux décimales $\og 60\fg{}$ sont les \emph{seuls} chiffres significatifs ; le $\og 3\fg{}$ ne fait
que situer l'ordre de grandeur ($10^{-4}$).
\end{exemple}

\begin{exemple}[sens inverse : de pH à concentration]
Une solution a un $\mathrm{pH}=2{,}00$ (2 décimales). On remonte à la concentration :
\[ [\ce{H3O+}]=10^{-\mathrm{pH}}=10^{-2,00}=\num{1,0e-2}\ \si{\mol\per\litre}. \]
Les $2$ décimales du pH donnent $2$ \CS\ à la concentration : on écrit $\num{1,0e-2}$ et non
$\num{1e-2}$.
\end{exemple}

% =====================================================================
\section{Justesse, fidélité et incertitude}
% =====================================================================
Les \CS\ traduisent la \emph{précision} d'écriture d'un nombre. Mais une mesure peut être
\emph{précise} sans être \emph{correcte}. Deux notions complémentaires permettent de qualifier la
qualité d'une mesure.

\begin{definition}[justesse et fidélité]
\begin{itemize}[leftmargin=1.6em,topsep=2pt,itemsep=2pt]
  \item la \textbf{justesse} (en anglais \emph{accuracy}) : proximité de la mesure à la \emph{vraie}
        valeur. Une mesure juste n'est pas \og biaisée\fg{} ;
  \item la \textbf{fidélité} (en anglais \emph{precision}) : proximité des mesures \emph{entre elles}
        quand on répète l'expérience. Une mesure fidèle est \og reproductible\fg{}.
\end{itemize}
\end{definition}

\begin{center}
\begin{tikzpicture}[scale=0.8,
   cible/.style={fill=primary},
   anneaux/.pic={
     \fill[red!12] (0,0) circle (1.5);
     \fill[white]  (0,0) circle (1.0);
     \fill[red!12] (0,0) circle (0.5);
     \draw[black!60] (0,0) circle (1.5);
     \draw[black!60] (0,0) circle (1.0);
     \draw[black!60] (0,0) circle (0.5);
     \fill[black!70] (0,0) circle (0.06);}]
  % --- 1 : juste & fidèle (points groupés au centre) ---
  \begin{scope}[xshift=0cm]
     \pic{anneaux};
     \fill[cible] (0,0) circle (0.09);
     \fill[cible] (0.18,0.12) circle (0.09);
     \fill[cible] (-0.12,0.16) circle (0.09);
     \fill[cible] (0.1,-0.15) circle (0.09);
     \node[below,align=center,font=\footnotesize] at (0,-1.95) {juste\\ \& fidèle};
  \end{scope}
  % --- 2 : fidèle mais non juste (points groupés, décentrés) ---
  \begin{scope}[xshift=4.2cm]
     \pic{anneaux};
     \fill[cible] (1.0,0.9) circle (0.09);
     \fill[cible] (1.18,1.05) circle (0.09);
     \fill[cible] (0.86,1.0) circle (0.09);
     \fill[cible] (1.05,0.78) circle (0.09);
     \node[below,align=center,font=\footnotesize] at (0,-1.95) {fidèle\\ mais non juste};
  \end{scope}
  % --- 3 : juste mais non fidèle (points dispersés, centrés en moyenne) ---
  \begin{scope}[xshift=8.4cm]
     \pic{anneaux};
     \fill[cible] (-0.9,0.3) circle (0.09);
     \fill[cible] (0.8,-0.6) circle (0.09);
     \fill[cible] (0.2,1.0) circle (0.09);
     \fill[cible] (-0.4,-0.9) circle (0.09);
     \node[below,align=center,font=\footnotesize] at (0,-1.95) {juste\\ mais non fidèle};
  \end{scope}
  % --- 4 : ni juste ni fidèle (dispersés et décentrés) ---
  \begin{scope}[xshift=12.6cm]
     \pic{anneaux};
     \fill[cible] (0.9,1.0) circle (0.09);
     \fill[cible] (1.3,0.6) circle (0.09);
     \fill[cible] (-1.1,-0.4) circle (0.09);
     \fill[cible] (0.4,-1.2) circle (0.09);
     \node[below,align=center,font=\footnotesize] at (0,-1.95) {ni juste\\ ni fidèle};
  \end{scope}
\end{tikzpicture}
\end{center}

\noindent Le \og centre de la cible\fg{} représente la vraie valeur ; chaque point est une mesure
répétée. \textbf{Fidélité} = points serrés ; \textbf{justesse} = points centrés.

\begin{remarque}
Le nombre de \CS\ est une indication de \textbf{fidélité} (jusqu'où on sait répéter la lecture), pas de
\textbf{justesse}. Une balance déréglée peut afficher $\qty{13{,}02}{\gram}$ ($4$ \CS, très \og précis\fg)
tout en étant \emph{fausse} de $\qty{2}{\gram}$. Les \CS\ ne corrigent jamais un biais.
\end{remarque}

\subsection{Incertitude absolue et relative}

\begin{definition}[incertitude absolue]
L'\textbf{incertitude absolue} $\Delta x$ est l'\og épaisseur de doute\fg{} sur la mesure ; le résultat
s'écrit $x=x_{\text{mes}}\pm\Delta x$, avec $\Delta x$ dans la même unité que $x$. Par convention, le
dernier \cs\ écrit fixe l'ordre de l'incertitude.
\end{definition}

\begin{exemple}
$m=\qty{13{,}02}{\gram}$ sous-entend $\Delta m=\qty{0,01}{\gram}$, donc
$m=\qty{13{,}02\pm0{,}01}{\gram}$ : la vraie masse est entre $\qty{13{,}01}{}$ et $\qty{13{,}03}{\gram}$.
\end{exemple}

\begin{formule}[incertitude relative]
\[ \boxed{\;\dfrac{\Delta x}{x}\;}\qquad\text{(sans unité, souvent en \%)} \]
\begin{ou}
  \item $\Delta x$ : incertitude absolue, même unité que $x$ ;
  \item $x$ : valeur mesurée ;
  \item $\dfrac{\Delta x}{x}$ : poids \emph{relatif} du doute, qui permet de comparer la qualité de
        mesures de grandeurs différentes.
\end{ou}
\end{formule}

\begin{exemple}[précision relative d'une pesée]
Pour $m=\qty{13{,}02\pm0{,}01}{\gram}$ :
\[ \frac{\Delta m}{m}=\frac{\num{0,01}}{\num{13,02}}=\num{7,68e-4}\approx \qty{0,077}{\percent}. \]
Une incertitude relative de l'ordre de $\qty{0,1}{\percent}$ correspond bien à une mesure à $4$ \CS :
\emph{plus il y a de \CS, plus l'incertitude relative est faible}. À l'inverse, une mesure à $2$ \CS\
a typiquement une incertitude relative de l'ordre de $1\%$.
\end{exemple}

\begin{theoreme}[lien \CS\ $\leftrightarrow$ incertitude relative]
À $n$ \CS, l'incertitude relative est de l'ordre de $10^{-(n-1)}$ :
\[ 2\ \text{CS}\sim 1\%,\qquad 3\ \text{CS}\sim 0{,}1\%,\qquad 4\ \text{CS}\sim 0{,}01\%. \]
Le nombre de \CS\ est donc une \emph{lecture directe} de la précision relative d'une mesure.
\end{theoreme}

% =====================================================================
\section{Synthèse}
% =====================================================================

\begin{tcolorbox}[boxbase,colback=primary!5,colframe=primary,
   title={\textsc{Mémo — compter les chiffres significatifs}}]
\begin{itemize}[leftmargin=1.6em,topsep=2pt,itemsep=2pt]
  \item \textbf{Tout chiffre non nul} compte \emph{toujours}.
  \item \textbf{Zéros de tête} (à gauche du 1\textsuperscript{er} non nul) : \emph{jamais} significatifs.
  \item \textbf{Zéros captifs} (entre deux non nuls) : \emph{toujours} significatifs.
  \item \textbf{Zéros de queue après la virgule} : \emph{toujours} significatifs.
  \item \textbf{Zéros de queue d'un entier} (sans virgule) : \emph{ambigus} $\to$ notation scientifique.
  \item \textbf{Nombres exacts} (comptage, conversions) : \emph{infinité} de \CS, ne limitent rien.
  \item \textbf{Le déplacement de la virgule} (changement d'unité) ne change \emph{jamais} le nombre de \CS.
\end{itemize}
\end{tcolorbox}

\begin{tcolorbox}[boxbase,colback=formulacol!6,colframe=formulacol,
   title={\textsc{Mémo — opérations}}]
\begin{center}\renewcommand{\arraystretch}{1.4}
\begin{tabular}{@{}>{\raggedright\arraybackslash}p{3.7cm} >{\raggedright\arraybackslash}p{3.5cm} >{\raggedright\arraybackslash}p{6.3cm}@{}}
\toprule
\textbf{Opération} & \textbf{On compte\dots} & \textbf{Le résultat a\dots}\\
\midrule
Addition / soustraction & les \emph{décimales} & le minimum de décimales des données\\
Multiplication / division & les \emph{\CS} & le minimum de \CS\ des données\\
Logarithme ($\mathrm{pH}$) & les \CS\ de l'argument & autant de \emph{décimales} qu'il y a de \CS\\
\bottomrule
\end{tabular}
\end{center}
\textbf{Toujours :} calculer sans arrondir, puis arrondir \emph{une seule fois} à la fin.
\end{tcolorbox}

\begin{tcolorbox}[boxbase,colback=defcol!5,colframe=defcol,
   title={\textsc{Mémo — vocabulaire à ne pas confondre}}]
\begin{itemize}[leftmargin=1.6em,topsep=2pt,itemsep=1.5pt]
  \item \textbf{Chiffre certain} vs \textbf{chiffre incertain} (le dernier écrit, estimé).
  \item \textbf{Chiffre significatif} vs \textbf{décimale} : on ne compte pas la même chose.
  \item \textbf{Justesse} (proche du vrai) vs \textbf{fidélité} (mesures groupées).
  \item \textbf{Incertitude absolue} $\Delta x$ (en unités) vs \textbf{relative} $\Delta x/x$ (en \%).
\end{itemize}
\end{tcolorbox}

\vfill
\begin{center}
\textcolor{primary}{\rule{0.5\textwidth}{1pt}}\\[4pt]
{\small\color{black!60}Fin de la Fiche 1 — Les chiffres significatifs.}
\end{center}

\end{document}
